% Standalone source. Compile with pdfLaTeX (or the Codex LaTeX preview).
% Evidence cutoff: 16 October 2026. No external assets required.
\documentclass[11pt,a4paper]{article}
\usepackage[T1]{fontenc}
\usepackage[utf8]{inputenc}
\usepackage{lmodern}
\usepackage[margin=18mm,headheight=14pt,footskip=10mm]{geometry}
\usepackage{array,booktabs,tabularx}
\usepackage[table]{xcolor}
\usepackage{microtype}
\usepackage{fancyhdr}
\usepackage{hyperref}
\definecolor{navy}{HTML}{183346}
\definecolor{teal}{HTML}{007E80}
\definecolor{pale}{HTML}{EDF5F5}
\definecolor{muted}{HTML}{566773}
\hypersetup{colorlinks=true,urlcolor=teal,linkcolor=teal,pdftitle={Learning to Use What a Model Knows},pdfauthor={Daniel Chen}}
\renewcommand{\familydefault}{\sfdefault}
\color{navy}
\setlength{\parindent}{0pt}
\setlength{\parskip}{6pt}
\setlength{\tabcolsep}{5pt}
\renewcommand{\arraystretch}{1.24}
\pagestyle{fancy}\fancyhf{}
\renewcommand{\headrulewidth}{0pt}
\renewcommand{\footrulewidth}{0.3pt}
\fancyfoot[L]{\scriptsize Daniel Chen | Compositional factual reasoning | 16 October 2026}
\fancyfoot[R]{\scriptsize\thepage}
\newcommand{\reportheading}[1]{\par\vspace{5pt}{\color{teal}\bfseries\large #1}\par\vspace{2pt}}
\newcommand{\reporttitle}[1]{{\fontsize{20}{24}\selectfont\bfseries #1}\par\vspace{8pt}}
\begin{document}
{\small\color{muted} RESEARCH SUMMARY / 01\par}
{\fontsize{28}{32}\selectfont\bfseries Learning to Use\newline What a Model Knows}\par\vspace{10pt}
\reportheading{Self-generated supervision for compositional factual reasoning}
{\small\color{muted} Daniel Chen | Research summary for faculty review\newline Evidence cutoff: 16 October 2026\par}
\reportheading{The question}
Can a language model improve how it combines facts using supervision generated by the same model, without external answer labels in its training targets? The project studies two-hop factual questions and follows the evidence from measuring failures to controlled training, transfer, and mechanism diagnosis.

\reportheading{The answer so far}
\textbf{There is a reproducible but modest benefit on the original task.} Context-conditioned teacher answers outperform direct self-distillation even after approximate matching of target accuracy. This is evidence for the value of how self-generated targets are constructed; it does not yet identify the internal mechanism or establish general multi-hop improvement.

{\small
\begin{tabular}{>{\raggedright\arraybackslash}p{\dimexpr 0.25458\linewidth-2\tabcolsep\relax}>{\raggedright\arraybackslash}p{\dimexpr 0.38289\linewidth-2\tabcolsep\relax}>{\raggedright\arraybackslash}p{\dimexpr 0.36253\linewidth-2\tabcolsep\relax}}
\rowcolor{navy}
\color{white}\bfseries Finding & \color{white}\bfseries Principal result & \color{white}\bfseries Interpretation \\
\rowcolor{pale}
Compositional failure persists & 33.80\% failure among 2,379 Qwen3-14B items with both hops correct & Operational knowledge does not guarantee composition. \\
Self-target training helps & C accuracy: 49.60\% to 53.13\%; +3.53 pp & Three-seed, within-dataset double holdout. \\
\rowcolor{pale}
Teacher context matters & Matched exp12: +2.48 pp over direct targets; 95\% interval [1.83, 3.12] & Six seeds; fixed 303-item reused evaluation subset. \\
Broad transfer is unconfirmed & MuSiQue two-hop EM: +0.217 pp; [-0.087, 0.521] & Predefined transfer criterion not met. \\
\end{tabular}
}\par\smallskip
{\small\color{muted} Sources: E1, E8-9, E12m, E13. pp = percentage points. Intervals here describe training-seed variation on fixed items, not population-wide uncertainty. Source key on page 9.\par}
\par\smallskip\noindent\colorbox{pale}{\parbox{\dimexpr\linewidth-2\fboxsep\relax}{\textbf{Research contribution at this stage:} an empirical account of when self-generated supervision helps, which simple explanations are weakened by controls, and where the effect fails to generalize. The broader claim of autonomous self-improvement remains open.}}\par\smallskip
{\small\color{muted} Reading guide: setup and measurement (p. 2); training controls (pp. 3-4); teacher-source comparison (p. 5); external transfer (p. 6); current diagnosis (p. 7); historical evidence and conclusions (p. 8); sources (p. 9).\par}

\clearpage
{\small\color{muted} RESEARCH SUMMARY / 02\par}
\reporttitle{1. From factual knowledge to composition}
\reportheading{Task and training method}
The original benchmark is Compositional Celebrities [1]. A question combines a person-to-birth-country fact with a country attribute, such as its capital. The compositionality gap is the proportion of composed questions answered incorrectly \textit{among items whose two constituent questions are both answered correctly}. ``Knowing'' is therefore an operational property under a fixed prompt and decoding setup.

\par\smallskip\noindent\colorbox{pale}{\parbox{\dimexpr\linewidth-2\fboxsep\relax}{\textbf{Self/context teacher:} answer hop 1 \(\rightarrow\) place the model's own answer in a background statement \(\rightarrow\) answer the composed question \(\rightarrow\) train a student to answer the original question without that background.\newline \textbf{Direct teacher control:} generate the target directly from the composed question. Both use the same base model; neither receives the gold intermediate fact when generating self targets.}}\par\smallskip
Training uses Qwen3-14B with LoRA rank 16, alpha 32, dropout 0.05, AdamW at 10\textsuperscript{-4}, effective batch 8, and two epochs per fit unless stated otherwise. Targets are short answers; prompt tokens are masked from the loss. Generation is greedy, capped at 24 new tokens, without explicit long chain-of-thought. The original full pool is 864 questions from 216 training people. [E8-12]

Training relation families: capital, calling code, and coordinates. Held-out families: currency, top-level domain, and numeric country code. A holds out people; B holds out relation families; C holds out both (500 questions each). Countries and country facts can overlap across splits. Exp12 adds a separately reserved D set of 500 question identities; this remains a within-dataset evaluation.

\reportheading{Exp1 and exp7: a real gap, with measurement caveats}
{\small
\begin{tabular}{>{\raggedright\arraybackslash}p{\dimexpr 0.23625\linewidth-2\tabcolsep\relax}>{\raggedright\arraybackslash}p{\dimexpr 0.30346\linewidth-2\tabcolsep\relax}>{\raggedright\arraybackslash}p{\dimexpr 0.22607\linewidth-2\tabcolsep\relax}>{\raggedright\arraybackslash}p{\dimexpr 0.23422\linewidth-2\tabcolsep\relax}}
\rowcolor{navy}
\color{white}\bfseries Model & \color{white}\bfseries Both hops correct (n) & \color{white}\bfseries Direct gap & \color{white}\bfseries Gold-context gap \\
\rowcolor{pale}
Qwen3-0.6B & 143 & 42.66\% & 60.14\% \\
Qwen3-1.7B & 947 & 74.66\% & 65.36\% \\
\rowcolor{pale}
Qwen3-8B & 1,902 & 39.54\% & 12.83\% \\
Qwen3-14B & 2,379 & 33.80\% & 10.93\% \\
\end{tabular}
}\par\smallskip
{\small\color{muted} Revised signed-floor scoring on the clean birthplace task pool (4,207 questions). Each row has a different known-item subset. Gold-context evaluation supplies the reference intermediate fact; this is distinct from the self-generated background used in training. [E1]\par}
\textbf{Fair scale comparison:} on the same 738 items known by 1.7B, 8B and 14B, direct gaps are 74.12\%, 26.83\% and 15.45\%. Scale clearly helps on this shared subset. A model-assisted audit of 120 sampled failures found 94 substantive errors, 12 reference ambiguities, 8 format cases and 6 scoring errors; it was not an independent blind human audit. [E1, E7]

\reportheading{Exp2: readable knowledge is not necessarily usable knowledge}
After excluding intermediate/final-answer first-token overlap, an intermediate entity's first token appears in a layer's top 10 for 243/288 failures (84.38\%) and 254/278 successes (91.37\%). This is a readout diagnostic, not proof of full-entity representation, causal use, or a unique ``too-late'' failure mechanism. [E2]


\clearpage
{\small\color{muted} RESEARCH SUMMARY / 03\par}
\reporttitle{2. What makes self-training useful?}
\reportheading{Exp6 \(\rightarrow\) exp8: from an initial signal to multi-seed evidence}
Early self, gold-label and hop-only training motivated the later controlled series. Those early contrasts remain historical evidence (p. 8); the main quantitative training claims below use verified exp8-10 records.

{\small
\begin{tabular}{>{\raggedright\arraybackslash}p{\dimexpr 0.36864\linewidth-2\tabcolsep\relax}>{\raggedright\arraybackslash}p{\dimexpr 0.14664\linewidth-2\tabcolsep\relax}>{\raggedright\arraybackslash}p{\dimexpr 0.31568\linewidth-2\tabcolsep\relax}>{\raggedright\arraybackslash}p{\dimexpr 0.16904\linewidth-2\tabcolsep\relax}}
\rowcolor{navy}
\color{white}\bfseries Evaluation split & \color{white}\bfseries Base & \color{white}\bfseries Full self, mean \(\pm\) SD & \color{white}\bfseries Gain \\
\rowcolor{pale}
A: new people & 39.00\% & 47.87\% \(\pm\) 0.50 & +8.87 pp \\
B: new relation families & 44.60\% & 50.80\% \(\pm\) 1.60 & +6.20 pp \\
\rowcolor{pale}
C: both new & 49.60\% & 53.13\% \(\pm\) 0.90 & +3.53 pp \\
\end{tabular}
}\par\smallskip
{\small\color{muted} Three training seeds; 500 items per split; 864 training examples. C seeds: 53.2\%, 52.2\%, 54.0\%; gain interval [1.29, 5.77] pp. Exp8 and exp9 full-self summaries are the same evidence, not two independent replications. \(\pm\) denotes sample standard deviation, not a confidence interval. [E8-9]\par}
\reportheading{Exp8: cleaner labels do not guarantee better transfer}
{\small
\begin{tabular}{>{\raggedright\arraybackslash}p{\dimexpr 0.51731\linewidth-2\tabcolsep\relax}>{\raggedright\arraybackslash}p{\dimexpr 0.16701\linewidth-2\tabcolsep\relax}>{\raggedright\arraybackslash}p{\dimexpr 0.31568\linewidth-2\tabcolsep\relax}}
\rowcolor{navy}
\color{white}\bfseries Target selection & \color{white}\bfseries Training n & \color{white}\bfseries C accuracy, mean \(\pm\) SD \\
\rowcolor{pale}
Full self & 864 & 53.13\% \(\pm\) 0.90 \\
Global confidence, top 50\% & 432 & 53.47\% \(\pm\) 0.95 \\
\rowcolor{pale}
Global random, 50\% & 432 & 54.07\% \(\pm\) 0.50 \\
Global confidence, top 70\% & 605 & 52.80\% \(\pm\) 0.72 \\
\rowcolor{pale}
Global random, 70\% & 605 & 54.53\% \(\pm\) 0.31 \\
\end{tabular}
}\par\smallskip
For seed 0, top-half confidence filtering raises target accuracy from 51.50\% to 70.60\%, but it does not outperform equally sized random selection on C. Filtering also changes category composition, motivating exp9. [E8]

\reportheading{Exp9: question-target correspondence matters}
Within-category confidence and random half-sampling give 52.53\% \(\pm\) 1.79 and 52.47\% \(\pm\) 0.90 on C. In contrast, shuffling full-pool targets within category collapses accuracy to 5.20\% \(\pm\) 1.51, versus 53.13\% for full self. Preserving the answer distribution alone is insufficient; the pairing matters. This severe mismatch control does not isolate a specific neural mechanism. [E9]

\reportheading{Exp10: even targets scored wrong can yield a small gain}
{\small
\begin{tabular}{>{\raggedright\arraybackslash}p{\dimexpr 0.40122\linewidth-2\tabcolsep\relax}>{\raggedright\arraybackslash}p{\dimexpr 0.14257\linewidth-2\tabcolsep\relax}>{\raggedright\arraybackslash}p{\dimexpr 0.24644\linewidth-2\tabcolsep\relax}>{\raggedright\arraybackslash}p{\dimexpr 0.20978\linewidth-2\tabcolsep\relax}}
\rowcolor{navy}
\color{white}\bfseries Category-matched condition & \color{white}\bfseries Training n & \color{white}\bfseries C accuracy & \color{white}\bfseries Gain over base \\
\rowcolor{pale}
Own wrong targets & 228 & 51.53\% \(\pm\) 0.42 & +1.93 pp \\
Own correct targets & 228 & 51.80\% \(\pm\) 0.69 & +2.20 pp \\
\rowcolor{pale}
Shuffled wrong targets & 228 & 8.73\% \(\pm\) 0.31 & -40.87 pp \\
\end{tabular}
}\par\smallskip
{\small\color{muted} Wrong minus correct: -0.27 pp, 95\% seed interval [-1.52, 0.98]. This is not an equivalence result. Correct/wrong groups match category counts but use different questions; gold labels are used to select these diagnostic subsets. Automatic ``wrong'' is not independently adjudicated factual wrongness. [E10]\par}

\clearpage
{\small\color{muted} RESEARCH SUMMARY / 04\par}
\reporttitle{3. Can the gain compound through iteration?}
\reportheading{Exp11: three routes, three seeds, a bounded checkpoint}
The same 864 questions are reused across rounds. \textbf{Reset} trains from the base model on refreshed targets; \textbf{continue} inherits the adapter but resets optimizer and schedule; \textbf{frozen} trains on initial targets for matched cumulative gradient steps. Shared round 1 plus six later conditions produces 21 fits. This compares target refresh against more training, while generation cost and optimizer schedules are not fully matched. [E11]

{\small
\begin{tabular}{>{\raggedright\arraybackslash}p{\dimexpr 0.44807\linewidth-2\tabcolsep\relax}>{\raggedright\arraybackslash}p{\dimexpr 0.31161\linewidth-2\tabcolsep\relax}>{\raggedright\arraybackslash}p{\dimexpr 0.24033\linewidth-2\tabcolsep\relax}}
\rowcolor{navy}
\color{white}\bfseries Condition & \color{white}\bfseries C accuracy, mean \(\pm\) SD & \color{white}\bfseries Hop-2 accuracy, mean \\
\rowcolor{pale}
Base & 50.00\% & 78.40\% \\
Shared round 1 & 52.53\% \(\pm\) 1.33 & 77.07\% \\
\rowcolor{pale}
Reset round 2 & 52.93\% \(\pm\) 1.10 & 76.67\% \\
Reset round 3 & 53.53\% \(\pm\) 0.50 & 76.60\% \\
\rowcolor{pale}
Continue round 2 & 53.67\% \(\pm\) 0.95 & 76.33\% \\
Continue round 3 & 53.67\% \(\pm\) 1.01 & 75.60\% \\
\rowcolor{pale}
Frozen targets, 2-round budget & 54.00\% \(\pm\) 1.06 & 77.53\% \\
Frozen targets, 3-round budget & 52.73\% \(\pm\) 1.36 & 75.00\% \\
\end{tabular}
}\par\smallskip
{\small\color{muted} All metrics use the same 500 C questions. The base is 50.0\%, not exp8's 49.6\%: regenerating base predictions changed 18 answer strings, with four wrong-to-right and two right-to-wrong flips. Gains must use each run's own baseline. [E11]\par}
\reportheading{The relevant comparison is against the budget control}
{\small
\begin{tabular}{>{\raggedright\arraybackslash}p{\dimexpr 0.56415\linewidth-2\tabcolsep\relax}>{\raggedright\arraybackslash}p{\dimexpr 0.18941\linewidth-2\tabcolsep\relax}>{\raggedright\arraybackslash}p{\dimexpr 0.24644\linewidth-2\tabcolsep\relax}}
\rowcolor{navy}
\color{white}\bfseries Paired C contrast & \color{white}\bfseries Difference & \color{white}\bfseries 95\% seed interval \\
\rowcolor{pale}
Reset R3 minus R1 & +1.00 pp & [-3.55, 5.55] \\
Reset R3 minus frozen 3-round & +0.80 pp & [-3.75, 5.35] \\
\rowcolor{pale}
Continue R3 minus frozen 3-round & +0.93 pp & [-3.05, 4.92] \\
\end{tabular}
}\par\smallskip
No test in the 54-comparison primary C family survives Holm correction at 0.05. The intervals are wide with three seeds; these results do not establish either an iterative benefit or equivalence. Frozen targets at the two-round budget have the highest observed C mean, but selecting this condition on the test results would be post hoc.

\reportheading{More training also has costs}
Hop-2 accuracy falls under both refreshed and frozen-target routes. Generated-target accuracy stays around 50-51\%, while the model becomes more confident in its scored-wrong training targets. This is compatible with self-reinforcement, but changing error sets prevent a simple calibration claim. The stage took about 5.34 hours across recorded generation, loading, training and evaluation phases. [E11]

\par\smallskip\noindent\colorbox{pale}{\parbox{\dimexpr\linewidth-2\fboxsep\relax}{\textbf{Checkpoint decision:} retain the modest one-stage benefit; do not claim that repeated target refresh produces reliable compounding improvement. The next question became why one source of self-generated targets is better than another.}}\par\smallskip

\clearpage
{\small\color{muted} RESEARCH SUMMARY / 05\par}
\reporttitle{4. Does the teacher's context make a difference?}
\reportheading{Exp12: direct self-distillation versus context-conditioned targets}
Both arms train on the same composed questions with paired LoRA initialization and matched update counts. Context uses the model's own first-hop answer, not gold context. The original three seeds were extended to six; then existing predictions were rescored and a separate matched-target training stage was run. These are distinct analyses, not interchangeable estimates. [E12a-c]

{\small
\begin{tabular}{>{\raggedright\arraybackslash}p{\dimexpr 0.36660\linewidth-2\tabcolsep\relax}>{\raggedright\arraybackslash}p{\dimexpr 0.13035\linewidth-2\tabcolsep\relax}>{\raggedright\arraybackslash}p{\dimexpr 0.13035\linewidth-2\tabcolsep\relax}>{\raggedright\arraybackslash}p{\dimexpr 0.37271\linewidth-2\tabcolsep\relax}}
\rowcolor{navy}
\color{white}\bfseries Stage / evaluation & \color{white}\bfseries Context & \color{white}\bfseries Direct & \color{white}\bfseries Difference; 95\% seed interval \\
\rowcolor{pale}
864 training items; D500; original scorer & 51.80\% & 49.27\% & +2.53 pp [0.71, 4.36] \\
Same predictions; D500; revised v3 scorer & 50.60\% & 45.70\% & +4.90 pp [3.28, 6.52] \\
\rowcolor{pale}
Same 864 training; fixed D303 subset; v3 & See contrast & See contrast & +3.14 pp [2.53, 3.75] \\
New matched785 training; fixed D303; v3 & 71.07\% & 68.59\% & +2.48 pp [1.83, 3.12] \\
\end{tabular}
}\par\smallskip
{\small\color{muted} All differences are context minus direct; unmatched D303 is reported as a contrast only. v3 is a post hoc, conservative whole-answer scorer; review-flagged outputs count as wrong in the main analysis. [E12b-c, E12m]\par}
\reportheading{Why the scoring audit matters}
The original scorer could accept substrings and incompatible answers. Under v3, 109-153 of the 500 outputs per adapted model are flagged for review. On the original six-seed predictions, jointly excluding either arm's review-flagged items gives +3.06 pp [2.46, 3.66], with 320-364 items per seed. The fixed 303-item subset excludes review flags across all 13 original models. It was selected after observing outputs, then frozen for matched785; it is not an untouched test set. [E12c]

\reportheading{Matched785: the strongest local control so far}
The two target sets have identical original-score correct counts (359/785 each); revised accepted counts are 360 versus 359. Each of 12 fits uses 785 items, two epochs and 198 updates. On fixed D303, the base scores 67.66\%; context gains +3.41 pp [3.05, 3.77] and direct +0.94 pp [0.53, 1.34]. All six paired context-minus-direct differences are positive (1.65-3.30 pp). [E12m]

Sensitivity checks remain positive: review flags counted correct yield +2.81 pp; excluding the paired union yields +2.60 pp. On full D500, the main difference is +4.30 pp [3.37, 5.23]. These scenarios are not bounds over every possible factual adjudication.

\par\smallskip\noindent\colorbox{pale}{\parbox{\dimexpr\linewidth-2\fboxsep\relax}{\textbf{What this weakens:} the simple explanation that context helps only because its targets have higher aggregate automatic accuracy.\newline \textbf{What remains:} 45 training items have review flags in at least one arm; target-token budgets differ by 6.27\%; semantic error types and conditional likelihood differ. Approximate score matching is not exact factual-quality matching, and the result does not prove a knowledge-retrieval mechanism.}}\par\smallskip

\clearpage
{\small\color{muted} RESEARCH SUMMARY / 06\par}
\reporttitle{5. Does the effect transfer to a new dataset?}
\reportheading{Exp13: a fixed external test, without further training}
The base model and all 12 matched785 adapters were evaluated on MuSiQue [4]: 1,229 two-hop, 751 three-hop and 405 four-hop questions. This experiment uses \textbf{closed-book questions without supporting passages}; its scores are not directly comparable to the official reading-comprehension setting. The primary endpoint was two-hop exact match (EM), using the frozen alias-aware scoring protocol. [E13]

{\small
\begin{tabular}{>{\raggedright\arraybackslash}p{\dimexpr 0.71487\linewidth-2\tabcolsep\relax}>{\raggedright\arraybackslash}p{\dimexpr 0.28513\linewidth-2\tabcolsep\relax}}
\rowcolor{navy}
\color{white}\bfseries Two-hop primary endpoint (n = 1,229) & \color{white}\bfseries EM \\
\rowcolor{pale}
Base model & 6.428\% \\
Context-trained adapters, six-seed mean & 6.564\% \\
\rowcolor{pale}
Direct-trained adapters, six-seed mean & 6.347\% \\
\end{tabular}
}\par\smallskip
{\small
\begin{tabular}{>{\raggedright\arraybackslash}p{\dimexpr 0.53564\linewidth-2\tabcolsep\relax}>{\raggedright\arraybackslash}p{\dimexpr 0.20367\linewidth-2\tabcolsep\relax}>{\raggedright\arraybackslash}p{\dimexpr 0.26069\linewidth-2\tabcolsep\relax}}
\rowcolor{navy}
\color{white}\bfseries Paired contrast & \color{white}\bfseries Difference & \color{white}\bfseries 95\% seed interval \\
\rowcolor{pale}
Context minus direct & +0.217 pp & [-0.087, 0.521] \\
Context minus base & +0.136 pp & [-0.116, 0.387] \\
\rowcolor{pale}
Direct minus base & -0.081 pp & [-0.243, 0.081] \\
\end{tabular}
}\par\smallskip
The predeclared success rule required both context-minus-direct and context-minus-base lower interval bounds to exceed zero. Neither does. The mean context advantage equals about 2.67 additional correct answers over direct and 1.67 over base per 1,229 items. This is a small point estimate with uncertainty, not proof of zero effect. [E13]

\reportheading{A small, predefined known-facts subset shows a signal}
Only 73/1,229 two-hop questions have both constituent answers correct under the base model. On these fixed items, EM is 24.66\% for base, 26.71\% for context and 24.43\% for direct: context minus direct is +2.28 pp [1.54, 3.03]. However, across all six seeds, the repaired baseline failures involve only \textbf{three distinct questions}. Repeating seeds does not create more independent examples.

The subset is informative about where improvement may occur, but cannot replace the full primary endpoint. Limited measured single-hop knowledge, question reformulation and answer-reference ambiguity may all contribute to the low overall scores; none cancels the observed primary result.

\reportheading{Secondary endpoints do not give one uniform story}
Context-minus-direct EM is -0.222 pp on three-hop questions and +0.123 pp on four-hop questions. Across all 2,385 questions, EM changes by +0.063 pp [-0.091, 0.217], while token F1 changes by +0.233 pp [0.122, 0.343]. These secondary analyses are not multiplicity-adjusted; exact correctness and token overlap should not be conflated. [E13]

\par\smallskip\noindent\colorbox{pale}{\parbox{\dimexpr\linewidth-2\fboxsep\relax}{\textbf{Decision:} preserve MuSiQue as a fixed, already-inspected transfer benchmark. The current evidence supports a local original-task effect, while overall cross-dataset transfer remains unsupported under the prespecified criterion. No tuning or rerunning of this benchmark was performed for this summary.}}\par\smallskip

\clearpage
{\small\color{muted} RESEARCH SUMMARY / 07\par}
\reporttitle{6. Current diagnosis: consistency is a hypothesis}
\reportheading{Why look beyond target correctness?}
Exp10 found small gains from targets scored wrong; exp12 found a context advantage after approximate target-accuracy matching. A plausible hypothesis is that a target can be factually wrong yet remain compatible with the model's own intermediate answer. An offline analysis examined existing outputs only; it did not train or query models again. [D16]

Of the matched 785 questions, 381 have identical target strings in both arms and 404 differ. On the 404 differing-target questions, context has 316 targets scored wrong (including review flags) and direct has 317. ``Consistency'' was tested against a dataset-derived country-attribute table and against saved model single-hop answers for the model's predicted country.

{\small
\begin{tabular}{>{\raggedright\arraybackslash}p{\dimexpr 0.35031\linewidth-2\tabcolsep\relax}>{\raggedright\arraybackslash}p{\dimexpr 0.24236\linewidth-2\tabcolsep\relax}>{\raggedright\arraybackslash}p{\dimexpr 0.24236\linewidth-2\tabcolsep\relax}>{\raggedright\arraybackslash}p{\dimexpr 0.16497\linewidth-2\tabcolsep\relax}}
\rowcolor{navy}
\color{white}\bfseries Consistency definition & \color{white}\bfseries Context wrong targets & \color{white}\bfseries Direct wrong targets & \color{white}\bfseries Difference \\
\rowcolor{pale}
Country-attribute table & 75/316 = 23.73\% & 6/317 = 1.89\% & +21.84 pp \\
Model single-hop, same run & 108/316 = 34.18\% & 16/317 = 5.05\% & +29.13 pp \\
\rowcolor{pale}
Model single-hop, older cache fills gaps & 142/316 = 44.94\% & 25/317 = 7.89\% & +37.05 pp \\
\end{tabular}
}\par\smallskip
{\small\color{muted} These are descriptive rates from one cached target pool, not six independent teacher samples. Unknown cases remain in the denominators. The primary attribute-table analysis has 36.71\% and 41.01\% unknown; arbitrary assignments to unknowns allow a difference from -19.17 to +58.55 pp. Older-cache evidence has better coverage but mixes runs. [D16]\par}
The two wrong-target sets differ. Restricting to the same 236 questions where both targets are wrong without a review flag gives table-compatible counts 42 versus 5, a +15.68 pp difference; same-run model consistency gives 74 versus 14, or +25.42 pp. These checks support a candidate association, not a mediation estimate of the downstream training gain.

\reportheading{What is completed, running, and only planned}
{\small
\begin{tabular}{>{\raggedright\arraybackslash}p{\dimexpr 0.30550\linewidth-2\tabcolsep\relax}>{\raggedright\arraybackslash}p{\dimexpr 0.36660\linewidth-2\tabcolsep\relax}>{\raggedright\arraybackslash}p{\dimexpr 0.32790\linewidth-2\tabcolsep\relax}}
\rowcolor{navy}
\color{white}\bfseries Workstream & \color{white}\bfseries Status at cutoff & \color{white}\bfseries How it enters the conclusions \\
\rowcolor{pale}
Exp14: Llama-3.1-8B & User reports run in progress; no completed results found locally & No cross-family replication claim yet. \\
Exp15 / exp16 & No separate result artifacts located under these identifiers & No missing numbers inferred; offline diagnosis is D16. \\
\rowcolor{pale}
Exp17: consistency intervention & Design, feasibility counts and data-preparation artifacts only & No causal outcome; no new HPC run launched here. \\
\end{tabular}
}\par\smallskip
The proposed exp17 control keeps the same target-string multiset and changes only its assignment to questions, under category and eligibility constraints. Existing same-run caches support 124 questions per arm under strict matching; larger-pool counts remain projections. This is intended to reduce source, length and format confounds before testing consistency causally. [P17]

\par\smallskip\noindent\colorbox{pale}{\parbox{\dimexpr\linewidth-2\fboxsep\relax}{\textbf{Current priority:} consolidate scoring/reference errors, factual coverage and composition failures; then make a bounded decision about one targeted repair. Preserve the frozen exp14 protocol and the already-inspected external benchmark.}}\par\smallskip

\clearpage
{\small\color{muted} RESEARCH SUMMARY / 08\par}
\reporttitle{7. Historical evidence and defensible conclusions}
\reportheading{Exp3-6b: retain the history, qualify the evidence}
The October 3 handoff records the following early results. The later exp1-11 report could not locate the corresponding original outputs for several interventions and early controls. They explain the research path, but are \textbf{not treated as reverified quantitative evidence} or used to establish a causal mechanism here. [H, R11]

{\small
\begin{tabular}{>{\raggedright\arraybackslash}p{\dimexpr 0.23014\linewidth-2\tabcolsep\relax}>{\raggedright\arraybackslash}p{\dimexpr 0.41955\linewidth-2\tabcolsep\relax}>{\raggedright\arraybackslash}p{\dimexpr 0.35031\linewidth-2\tabcolsep\relax}}
\rowcolor{navy}
\color{white}\bfseries Experiment & \color{white}\bfseries Historical result & \color{white}\bfseries Evidence limit \\
\rowcolor{pale}
Exp3: per-item best-layer backpatching & 37\% repair versus 24\% control & Exploratory layer selection can inflate effects. \\
Exp4: fixed-layer backpatching & 14-16\% repair versus 3\% control; 5-9\% damage to correct items & Reported negative net utility; denominators and causal inference need raw-data verification. \\
\rowcolor{pale}
Exp5: confidence-gated patching & AUROC about 0.75; net gain about zero & Historical summary only. \\
Exp6: self / gold / hop-only & C base 49.6\%; self 54.2\% and 52.2\%; gold 56.8\%; hop-only 49.8\% & Early seeds and controls are not substituted for later verified runs. \\
\rowcolor{pale}
Exp6b: post-training lens & Self reportedly shifts intermediate readout slightly later; gold earlier & Does not establish that training repairs a timing bottleneck. \\
\end{tabular}
}\par\smallskip
\reportheading{What the completed evidence supports}
\textbf{1. A measurable composition problem.} Conditional failures remain after revised scoring. Larger models and supplied intermediate facts reduce the gap in shared-item comparisons; the data do not support ``scale cannot help.''

\textbf{2. A limited training effect.} Self-generated short-answer supervision improves original-task held-out scores. Context-conditioned targets outperform direct targets in six paired seeds, including an approximately accuracy-matched stage. The strongest claim is about this training procedure and evaluation setting.

\textbf{3. Useful negative and limiting results.} Confidence filtering and iterative refreshing do not show stable extra benefit. Longer training can reduce single-hop performance. The predefined overall MuSiQue transfer criterion is not met.

\textbf{4. An unresolved explanation.} Intermediate-token readability and target consistency are clues. Neither identifies a unique root cause. The effect could reflect several interacting factors: knowledge coverage, reference conventions, target semantics, confidence, and question-answer compatibility.

\reportheading{What a preprint should and should not claim}
A defensible framing is \textit{a controlled empirical study of self-generated supervision for compositional factual QA, including its gains and limits}. It should not claim a novel backpatching method, general autonomous self-improvement, a proven ``hopping too late'' mechanism, equal factual quality between matched targets, or demonstrated broad transfer.

{\small\color{muted} Before manuscript submission, the key evidence gaps are independent factual adjudication of disputed examples, completion and frozen-protocol analysis of exp14, and a clear distinction between prospective tests and post hoc analyses. The report makes no claim that these unfinished items are already resolved.\par}

\clearpage
{\small\color{muted} RESEARCH SUMMARY / 09\par}
\reporttitle{Evidence record and selected references}
\reportheading{How to read the statistics}
{\small\color{muted} All accuracies are percentages; differences are percentage points. Unless explicitly stated, \(\pm\) is sample SD across training seeds. Reported 95\% intervals are paired t intervals over three or six seeds on fixed items. They do not capture all item sampling, person/country clustering, target construction, reference truth, or dataset-selection uncertainty. A non-significant contrast is not evidence of equivalence. Repeated items and seeds are not independent new facts.\par}
\reportheading{Source key and verification scope}
{\small
\begin{tabular}{>{\raggedright\arraybackslash}p{\dimexpr 0.17923\linewidth-2\tabcolsep\relax}>{\raggedright\arraybackslash}p{\dimexpr 0.82077\linewidth-2\tabcolsep\relax}}
\rowcolor{navy}
\color{white}\bfseries Code & \color{white}\bfseries Local evidence used \\
\rowcolor{pale}
E1 / E7 & Revised gap comparison, common-item analysis and 120-item label audit; consolidated in the verified 11 October report. \\
E2 & Original lens items plus saved overlap-filtered analysis; filtering was not rebuilt with a fresh tokenizer for this report. \\
\rowcolor{pale}
E8 / E9 / E10 & Saved exp8 summary, exp9 verified report, exp10 verified report; the full-self overlap is counted only once. \\
E11 / R11 & Exp11 archived analysis and the 11 October consolidated report; includes raw-output rescoring verification and baseline reconciliation. \\
\rowcolor{pale}
E12a / E12b & Job700454 review and job700972 seed extension; original scorer and six-seed outputs. \\
E12c & Job701713 v3 rescoring and offline review-flag sensitivity analysis; post hoc scoring revision preserved alongside original scores. \\
\rowcolor{pale}
E12m & Job701792 matched785 local verification, independent analysis and target-label audit. \\
E13 & Job701955 local verification: frozen MuSiQue data, scoring and adapter receipts. \\
\rowcolor{pale}
D16 / P17 & 16 October offline consistency report; exp17 feasibility/design and preparation artifacts. \\
H & INWARD\_HANDOFF, dated 3 October; historical numbers only where raw verification is missing. \\
\end{tabular}
}\par\smallskip
{\small\color{muted} This summary synthesizes saved analyses and their recorded checks; it does not claim a new independent audit of every raw prediction. Later verification notes report archive/code/data hashes and independent recomputation of scores and seed intervals. Recent downloads omit adapter weights and final scheduler accounting; completion relies on program markers and logs. Exact local paths and SHA-256 hashes for the inputs used here are preserved in the companion source\_manifest.json. No original experimental outputs were overwritten, and no HPC experiments were launched.\par}
\reportheading{Position in the literature}
{\small\color{muted} The gap measurement follows [1]; internal intervention is related to [2]. The self-target method is closely related to context distillation [3], applied here to internally generated intermediate facts and short-answer targets. MuSiQue [4] supplies the external composition benchmark. These connections motivate the study; this brief is not an exhaustive novelty review.\par}
{\small\color{muted} \href{https://aclanthology.org/2023.findings-emnlp.378/}{[1] Press, O. et al. (2023). Measuring and Narrowing the Compositionality Gap in Language Models. Findings of EMNLP.}\par}
{\small\color{muted} \href{https://arxiv.org/abs/2406.12775}{[2] Biran, E. et al. (2024). Hopping Too Late: Exploring the Limitations of Large Language Models on Multi-Hop Queries. EMNLP.}\par}
{\small\color{muted} \href{https://arxiv.org/abs/2209.15189}{[3] Snell, C., Klein, D. and Zhong, R. (2022). Learning by Distilling Context.}\par}
{\small\color{muted} \href{https://aclanthology.org/2022.tacl-1.31/}{[4] Trivedi, H. et al. (2022). MuSiQue: Multihop Questions via Single-hop Question Composition. TACL 10:539-554.}\par}
\end{document}
